% Options for packages loaded elsewhere
\PassOptionsToPackage{unicode}{hyperref}
\PassOptionsToPackage{hyphens}{url}
\PassOptionsToPackage{dvipsnames,svgnames,x11names}{xcolor}
%
\documentclass[
]{article}
\usepackage{amsmath,amssymb}
\usepackage{iftex}
\ifPDFTeX
  \usepackage[T1]{fontenc}
  \usepackage[utf8]{inputenc}
  \usepackage{textcomp} % provide euro and other symbols
\else % if luatex or xetex
  \usepackage{unicode-math} % this also loads fontspec
  \defaultfontfeatures{Scale=MatchLowercase}
  \defaultfontfeatures[\rmfamily]{Ligatures=TeX,Scale=1}
\fi
\usepackage{lmodern}
\ifPDFTeX\else
  % xetex/luatex font selection
  \setmainfont[]{DejaVu Serif}
  \setmonofont[]{DejaVu Sans Mono}
\fi
% Use upquote if available, for straight quotes in verbatim environments
\IfFileExists{upquote.sty}{\usepackage{upquote}}{}
\IfFileExists{microtype.sty}{% use microtype if available
  \usepackage[]{microtype}
  \UseMicrotypeSet[protrusion]{basicmath} % disable protrusion for tt fonts
}{}
\makeatletter
\@ifundefined{KOMAClassName}{% if non-KOMA class
  \IfFileExists{parskip.sty}{%
    \usepackage{parskip}
  }{% else
    \setlength{\parindent}{0pt}
    \setlength{\parskip}{6pt plus 2pt minus 1pt}}
}{% if KOMA class
  \KOMAoptions{parskip=half}}
\makeatother
\usepackage{xcolor}
\usepackage[margin=1in]{geometry}
\setlength{\emergencystretch}{3em} % prevent overfull lines
\providecommand{\tightlist}{%
  \setlength{\itemsep}{0pt}\setlength{\parskip}{0pt}}
\setcounter{secnumdepth}{-\maxdimen} % remove section numbering
\usepackage{newunicodechar}
\newunicodechar{★}{\ensuremath{\star}}
\newunicodechar{✓}{\ensuremath{\checkmark}}
\newunicodechar{√}{\ensuremath{\surd}}
\newunicodechar{≈}{\ensuremath{\approx}}
\ifLuaTeX
  \usepackage{selnolig}  % disable illegal ligatures
\fi
\IfFileExists{bookmark.sty}{\usepackage{bookmark}}{\usepackage{hyperref}}
\IfFileExists{xurl.sty}{\usepackage{xurl}}{} % add URL line breaks if available
\urlstyle{same}
\hypersetup{
  colorlinks=true,
  linkcolor={Maroon},
  filecolor={Maroon},
  citecolor={Blue},
  urlcolor={Blue},
  pdfcreator={LaTeX via pandoc}}

\author{}
\date{}

\begin{document}

\hypertarget{chapter-1-scaled-dot-product-attention}{%
\section{Chapter 1 --- Scaled Dot-Product
Attention}\label{chapter-1-scaled-dot-product-attention}}

\begin{verbatim}
Type: theory + notebook
Ordering: theory-first
Depends on: 0 (conventions)
\end{verbatim}

\begin{quote}
\emph{``Attention is a differentiable, content-based lookup.''}
\end{quote}

\hypertarget{motivation}{%
\subsection{1.1 Motivation}\label{motivation}}

Given a sequence of token representations
\(\mathbf{x}_1, \dots, \mathbf{x}_n\), we want each token to reach into
the sequence and pull information from other tokens it finds relevant
--- based on \emph{content}, not position. Scaled dot-product attention
is the canonical solution: a differentiable, permutation-equivariant
operator that computes each output as a convex combination of value
vectors, weighted by query-key similarity.

\hypertarget{storyline}{%
\subsection{1.2 Storyline}\label{storyline}}

Attention is a four-step pipeline. Given queries
\(Q \in \mathbb{R}^{n \times d_k}\), keys
\(K \in \mathbb{R}^{m \times d_k}\), and values
\(V \in \mathbb{R}^{m \times d_v}\), each output \(\mathbf{o}_i\) is
produced by:

\begin{enumerate}
\def\labelenumi{\arabic{enumi}.}
\tightlist
\item
  \textbf{Score.} Compute
  \(s_{ij} = \langle \mathbf{q}_i, \mathbf{k}_j \rangle\) --- the raw
  content-match between query \(i\) and key \(j\). The matrix
  \(QK^\top \in \mathbb{R}^{n \times m}\) collects all scores.
\item
  \textbf{Scale.} Divide by \(\sqrt{d_k}\). This is the variance
  controller studied in Chapter 3 --- for now treat it as a fixed
  pre-softmax temperature.
\item
  \textbf{Normalize.} Apply row-wise softmax. The row
  \(\boldsymbol{\alpha}_i\) is now a probability distribution over the
  \(m\) source tokens --- the \emph{attention weights}.
\item
  \textbf{Mix.} Output
  \(\mathbf{o}_i = \sum_j \alpha_{ij} \mathbf{v}_j\) --- a convex
  combination of value rows, weighted by attention.
\end{enumerate}

The whole operation is \(\text{softmax}(QK^\top / \sqrt{d_k})\, V\).
Everything else in the transformer literature --- multi-head (Ch. 5),
causal masking (Ch. 4), positional encodings (Ch. 6), cross-attention,
KV-cache --- is a dressing of these four steps.

\textbf{Why this particular design.} Three forces are being balanced:
(a) \emph{content-addressing} (each query decides what to retrieve based
on its own state, not its position); (b) \emph{differentiability} (no
argmax --- the operation must admit gradients so the projection weights
can be learned); (c) \emph{permutation equivariance over sources} (the
operator is invariant to reordering keys/values as long as they are
permuted together, which forces positional information to be injected at
the input, not in the mechanism). These three pin down the
softmax-weighted convex combination uniquely up to parameterization.

\textbf{Train / val / inference.} The operator itself is identical
across the three regimes --- same four steps, same weights (once
trained). Two important wrinkles: - \textbf{Attention dropout} (zeroing
entries of \(\boldsymbol{\alpha}\), rescaling the rest) is applied
during training only. Validation and inference see the undrop'd weights.
- \textbf{KV-cache at inference.} During autoregressive generation,
\(K\) and \(V\) for past tokens are cached and re-used across steps;
only the new query interacts with the full key/value stack. This is an
engineering optimization that exploits the causal structure of Ch. 4,
not a change to the operator. - \textbf{Parallel vs.~sequential.}
Training runs attention on the full sequence in one shot; inference runs
one query per token sequentially (for autoregressive models). The
computation is identical, only the batching differs.

\hypertarget{rigorous-core}{%
\subsection{1.3 Rigorous core}\label{rigorous-core}}

\textbf{Definition 1.1 (Scaled dot-product attention).} Let
\(Q \in \mathbb{R}^{n \times d_k}\),
\(K \in \mathbb{R}^{m \times d_k}\),
\(V \in \mathbb{R}^{m \times d_v}\). Define
\[\text{Attention}(Q, K, V) = \text{softmax}\!\left( \frac{QK^\top}{\sqrt{d_k}} \right) V \in \mathbb{R}^{n \times d_v},\]
where softmax is applied row-wise on the \(m\)-dimensional rows of the
\(n \times m\) score matrix.

In self-attention, \(Q, K, V\) come from the same input
\(X \in \mathbb{R}^{n \times d}\) via learned projections
\(Q = X W_Q, K = X W_K, V = X W_V\) with
\(W_Q, W_K \in \mathbb{R}^{d \times d_k}\) and
\(W_V \in \mathbb{R}^{d \times d_v}\).

\textbf{Proposition 1.2 (Permutation equivariance over sources).} For
any \(m \times m\) permutation matrix \(P\),
\[\text{Attention}(Q, PK, PV) = \text{Attention}(Q, K, V).\]
\emph{Sketch: row-permutation equivariance of softmax plus
\(P^\top P = I\).} (Direct verification.)

\textbf{Corollary 1.2.1.} Self-attention is permutation-equivariant
under token reordering. Positional information must therefore be
injected \emph{before} attention (Chapter 6).

\textbf{Proposition 1.3 (Convex-hull containment).} Each output row
\(\mathbf{o}_i\) lies in the convex hull of
\(\{\mathbf{v}_1, \dots, \mathbf{v}_m\}\).

\textbf{Corollary 1.3.1.} For every coordinate \(k\),
\(\min_j V_{jk} \le (\mathbf{o}_i)_k \le \max_j V_{jk}\). Attention
cannot extrapolate beyond the value range --- a sharp, often overlooked
constraint.

\textbf{Proposition 1.4 (Smoothness).} \(\text{Attention}(Q, K, V)\) is
\(C^\infty\) in \(Q, K, V\), hence admits gradients everywhere.
(Composition of smooth maps; the failure of this property for hard
attention is what makes Exercise 1.4 interesting.)

\hypertarget{numerical-example-d_k-d_v-2-n-m-2-by-hand}{%
\subsection{\texorpdfstring{1.4 Numerical example ---
\(d_k = d_v = 2, n = m = 2\) by
hand}{1.4 Numerical example --- d\_k = d\_v = 2, n = m = 2 by hand}}\label{numerical-example-d_k-d_v-2-n-m-2-by-hand}}

Let
\[Q = \begin{pmatrix} 1 & 0 \\ 0 & 1 \end{pmatrix}, \quad K = \begin{pmatrix} 1 & 0 \\ 1 & 1 \end{pmatrix}, \quad V = \begin{pmatrix} 10 & 0 \\ 0 & 10 \end{pmatrix}.\]

\textbf{Step 1 --- score.}
\(QK^\top = \begin{pmatrix} 1 & 1 \\ 0 & 1 \end{pmatrix}\).

\textbf{Step 2 --- scale.} Divide by \(\sqrt{2} \approx 1.414\):
\(\begin{pmatrix} 0.707 & 0.707 \\ 0 & 0.707 \end{pmatrix}\).

\textbf{Step 3 --- softmax row-wise.} - Row 1: both entries equal, so
softmax gives \((0.5, 0.5)\). - Row 2:
\((e^0, e^{0.707}) / Z = (1, 2.028) / 3.028 \approx (0.330, 0.670)\).

\textbf{Step 4 --- mix with \(V\).} -
\(\mathbf{o}_1 = 0.5 \cdot (10, 0) + 0.5 \cdot (0, 10) = (5, 5)\). -
\(\mathbf{o}_2 = 0.330 \cdot (10, 0) + 0.670 \cdot (0, 10) \approx (3.30, 6.70)\).

\textbf{Sanity check against Corollary 1.3.1.} Each output coordinate is
in \([0, 10]\) --- the range of \(V\). ✓

\hypertarget{mechanics-check}{%
\subsection{1.5 Mechanics check}\label{mechanics-check}}

\begin{itemize}
\tightlist
\item
  \textbf{Score (\(QK^\top\)).} Measures content similarity between each
  query and each key via inner product. The Gram-matrix structure is
  what makes attention ``content-based'' rather than position-based.
\item
  \textbf{Scale (\(/\sqrt{d_k}\)).} Controls pre-softmax variance;
  without it, softmax saturates at large \(d_k\) (Ch. 3). No effect when
  \(d_k\) is small.
\item
  \textbf{Softmax.} Converts scores to a probability distribution.
  Enforces non-negativity, sum-to-one, differentiability, and
  order-preservation (Ch. 2). The \emph{only} step that is non-linear in
  \(Q, K\).
\item
  \textbf{Mix (\(\cdot\, V\)).} Produces output as a convex combination
  of value rows. By Corollary 1.3.1, the mechanism is a
  \emph{contraction} in the sense that outputs cannot exceed the value
  range.
\item
  \textbf{Separate projections \(W_Q, W_K, W_V\).} Decouples three
  roles: what a token asks for (Q), what it advertises as a matching
  index (K), what it contributes as payload (V). Tying any two collapses
  expressivity (Exercise 1.4).
\item
  \textbf{Row-wise softmax (not column-wise).} Normalizes across
  \emph{keys} per query, not across queries per key. Swapping this would
  break the ``each query pulls from the sequence'' semantics; the result
  would not even have the right units.
\end{itemize}

\hypertarget{exercises}{%
\subsection{1.6 Exercises}\label{exercises}}

\textbf{1.1 ★★} \emph{{[}Mechanism dissection --- hard-lookup limit.{]}}
Fix \(Q, K, V\) such that for some specific \(j^*\),
\(s_{ij^*} = \Delta\) and \(s_{ij} = 0\) for \(j \ne j^*\). Compute
\(\lim_{\Delta \to \infty} \mathbf{o}_i\). Now scale the scores by
\(\Delta\) itself: replace \(s_{ij}\) by \(\Delta \cdot s_{ij}\)
(equivalently, scale \(Q\) by \(\sqrt{\Delta}\)). Show that the output
converges to \(\mathbf{v}_{j^*}\) and quantify the rate: for \(n\) keys
with a margin-1 winner, how large must \(\Delta\) be to get
\(\alpha_{ij^*} > 1 - \epsilon\)? This is soft-argmax as the
temperature-zero limit, and the result is the bridge to Chapter 2.

\textbf{1.2 ★★} \emph{{[}Failure-mode construction --- tied
\(W_Q = W_K\).{]}} Suppose \(W_Q = W_K\) (tied), so
\(\mathbf{q}_i = \mathbf{k}_i\) for every \(i\). Show that the score
matrix is symmetric: \(s_{ij} = s_{ji}\). Construct a concrete \(n = 3\)
task for which untied \(W_Q, W_K\) is strictly more expressive --- i.e.,
the set of attention matrices realizable with tied projections is a
proper subset of those realizable with untied projections. What
structural property of \(QK^\top\) is being given up when we tie?
\emph{Hint: think about rank and asymmetry.}

\textbf{1.3 ★★★} \emph{{[}Theorem about the operator --- universality
and limits.{]}} Consider self-attention with fixed \(W_Q, W_K, W_V\) as
a map
\(X \in \mathbb{R}^{n \times d} \to Y \in \mathbb{R}^{n \times d_v}\).

\begin{enumerate}
\def\labelenumi{(\alph{enumi})}
\item
  Prove that \(Y\) always lies in the convex hull of
  \(\{W_V \mathbf{x}_1, \dots, W_V \mathbf{x}_n\}\), row-wise.
  (Corollary 1.3.1.)
\item
  Conclude that a single self-attention layer \emph{cannot} compute a
  function whose output range is outside the range of the linear map
  \(W_V\) applied to the input tokens. Give a concrete example of a
  function from \(\mathbb{R}^{n \times 1} \to \mathbb{R}^{n \times 1}\)
  that is representable by an MLP but not by any single self-attention
  layer.
\item
  What does this say about why transformers need both attention
  \emph{and} a feedforward block per layer? (You are previewing Chapter
  8.)
\end{enumerate}

\textbf{1.4 ★★★} \emph{{[}Predict-then-verify --- soft vs.~hard
attention.{]}} The notebook will compare three attention variants on a
simple retrieval task: (i) standard softmax attention, (ii) softmax with
temperature \(T \to 0\) (near-hard attention), and (iii) actual argmax
(non-differentiable). The task is: given \(n\) keys, one of which
matches the query, retrieve the corresponding value.

Predict, before running:

\begin{enumerate}
\def\labelenumi{(\alph{enumi})}
\tightlist
\item
  Which of (i)--(iii) reaches the best retrieval accuracy at
  convergence, and why?
\item
  Which trains fastest (fewest steps to near-zero loss)? Why?
\item
  Which \emph{fails to train at all}? Justify from Proposition 1.4 or
  its failure.
\item
  For small \(n\) (say \(n = 4\)), does (ii) strictly dominate (i)? What
  about large \(n\) (say \(n = 1024\))? Give a mechanistic reason for
  the \(n\)-dependence.
\end{enumerate}

State confidence per part (low / medium / high) with one-line reasons.
The notebook will test (a)--(d) on a toy copy task.

\hypertarget{before-the-notebook-01_attention_by_hand.ipynb}{%
\subsection{\texorpdfstring{1.7 Before the notebook
(\texttt{01\_attention\_by\_hand.ipynb})}{1.7 Before the notebook (01\_attention\_by\_hand.ipynb)}}\label{before-the-notebook-01_attention_by_hand.ipynb}}

Write down these predictions before opening the notebook:

\begin{enumerate}
\def\labelenumi{\arabic{enumi}.}
\tightlist
\item
  \textbf{Worked-example reproduction.} The numpy implementation should
  reproduce \((5, 5)\) and \(\approx (3.30, 6.70)\) from §1.4 to
  floating-point precision. Predict the specific floating-point error
  you expect (order of magnitude).
\item
  \textbf{Perfect-match retrieval.} With \(n = 4\) orthogonal keys and a
  query equal to key 2 scaled by \(\lambda\), predict the attention
  weights as \(\lambda\) varies over \(\{0.1, 1, 10, 100\}\).
  Specifically: at which \(\lambda\) does the attention become
  \(> 99\%\) concentrated on key 2?
\item
  \textbf{Equal keys.} If all keys are equal (so \(\mathbf{k}_j\) is the
  same vector for every \(j\)), predict the attention distribution and
  the output, without computation. Justify from the definition.
\item
  \textbf{Gradient sanity.} Prop. 2.2 gives
  \(\partial p_i / \partial s_j = p_i(\delta_{ij} - p_j)\). Predict
  which entries of the autograd-computed Jacobian will be exactly zero,
  positive, or negative, and check against PyTorch.
\item
  \textbf{Handcrafted self-attention.} The notebook sets up a 5-token
  toy sequence where tokens 0 and 3 share a signal. Predict which
  attention rows will concentrate mass on which source tokens
  \emph{before} running the cell.
\end{enumerate}

For each: direction, expected magnitude, and confidence level.

\end{document}
